\cotefile{29}
\scanpage{157}
\uline{Proposition 3.1.} (\unsure{Remontes|Remontées|Réunions}) $X$ étale sur $Y$ (par ex. fini).
$G$ groupe fini \deleted{\unsure{de|dé|dans}} \inserted{opérant par} $Y$-automorphismes sur $X$.

\deleted{Si $s\in$} \inserted{(i)} L'ens des pts de $X$ où $G$ opère
\inserted{\unsure{ou s'il y a lieu|où} \deleted{\unsure{sans}} \unsure{sans pts fixes}}
% slanted
\unsure{sans} ramification est à la fois ouvert et fermé.

(ii) \deleted{\unsure{$X/G$ et $X$}} $X/G$ existe, et $X\to X/G$
et $X/G\mathrel{\unsure{-|\to}}Y$ sont \uline{étales}.

Si $s\in G$ \deleted{\unsure{qui fixe}} est dans le groupe
d'inertie d'un pt $x\in X$, il \deleted{laisse fixe} \inserted{opère \unsure{triv.|laisse} sur \unsure{les|un pt}} \unsure{la}
comp. \unsure{str.|strte} \unsure{connexe|connexes} de $x$ (\unsure{cf|sic|ii}), donc l'ens $X(s)$
des $x\in X$ tels que $s\in G_{i}(x)$ est une
réunion de composantes connexes, donc ouvert et
fermé. Comme \deleted{\unsure{$X(s)$}} l'ens des pts où
$G_i(x)\ne e$ est la réunion des $X(s)$ $(s\ne e)$,
(i) s'ensuit. D'autre part, $X$ \unsure{étant|etant}
\unsure{quasi-affine|qu. affine} sur $Y$, $X/G$ existe $[\,\unsure{\text{---}}\,]$
\deleted{\unsure{donc}} \unsure{on va prouver|on se borne au|en faisant} \unsure{ce qui était|cas fini|ce que} \unsure{.}

Soit $\uline{I}$ l'ens des comp. connexes de $X$.
$G$ opère sur $\uline{I}$, \deleted{\unsure{et est réuni au cas|et $X$ est réunion des}} \inserted{et $X$ est réunion des}
\unsure{des} \deleted{\unsure{$X(\alpha)$|$X_{(\alpha)}$}} $\alpha\in I/G$, obtenus en
\deleted{\unsure{par.|posant} $\unsure{\alpha}$} $X(\alpha)=\bigcup_{i\in\alpha}X_{i}$, stable sous $G$.
donc $X/G=\bigcup X(\alpha)/G$ on est
% archive sheet number: 156
